\documentclass[12pt]{article}
\usepackage{amsmath}
\usepackage{amssymb}
\usepackage{physics}
\usepackage{hyperref}
\usepackage{graphicx}

% Shorthand for displaystyle fractions inside inline math
\newcommand{\ds}{\displaystyle}

\title{The Twin Paradox}
\author{}
\date{}

\begin{document}
\maketitle
% ============================================================
\section{Introduction}
% ============================================================

In this essay we will discuss the twin paradox, which arises in relativity when two objects,
initially at the same position and each equipped with its own clock, move away from each other
and later return to the same position. We assume that the two clocks are initially synchronized
and therefore measure exactly the same time. When the two objects meet again, however, the clocks
will in general have measured different elapsed times. \par

In the twin experiment, one of two identical twins embarks on a space journey, moving at a speed
close to the speed of light, while the other remains at home and waits for the traveling twin to
return. When the twins meet again, the twin who remained at home has aged more than the one who
traveled. This result appears paradoxical in two different ways. First, it is very counterintuitive
to accept that simply moving differently can change the amount of time experienced by the two twins.
Second, and more importantly, there appears to be a breaking of symmetry. Why should one twin age
differently from the other when there is no privileged system of reference? It seems as though the
younger twin must somehow be different from the twin who remained at home. Yet, during their relative
motion, each twin may regard the other as moving. We will examine in detail how this apparent
asymmetry arises and why there is no contradiction with the principle of relativity. \par

Below we will make use of formulas derived in the essay
\href{/physics/relativity/special-relativity/lorentz-transformations}{Lorentz Transformations}.
In particular, we will need the transformation between two inertial reference frames $F$ and $F'$,
where $F'$ moves with velocity $\vb{v}=v\vb{e}_1$ relative to $F$. The Lorentz transformation
for this boost is

\begin{equation}
  \begin{aligned}
    \boldsymbol{\Lambda}(\vb{v}) =\;
      &\gamma\ket{\vb{e}_1,0}\bra{\vb{e}_1,0}
      - \gamma\,\frac{v}{c}
        \ket{\vb{e}_1,0}\bra{\vb{0},1}
      + \ket{\vb{e}_2,0}\bra{\vb{e}_2,0}
      + \ket{\vb{e}_3,0}\bra{\vb{e}_3,0}
      \\
      &- \gamma\,\frac{v}{c}
        \ket{\vb{0},1}\bra{\vb{e}_1,0}
      + \gamma\ket{\vb{0},1}\bra{\vb{0},1}.
  \end{aligned}
  \label{eq:Lorentz_transformation}
\end{equation}

Here,

\begin{equation}
  \gamma
  =
  \frac{1}{\sqrt{1-\dfrac{v^2}{c^2}}}.
  \label{eq:Lorentz_factor}
\end{equation}

We will restrict our analysis to the case in which the relative velocity is directed along the
$x$-axis. In this case, the corresponding transformations of space and time are

\begin{equation}
  \begin{aligned}
    x' &= \gamma(x-vt), \\
    y' &= y, \\
    z' &= z, \\
    t' &= \gamma\!\left(t-\frac{v}{c^2}x\right).
  \end{aligned}
  \label{eq:Lorentz_transformation_along_x}
\end{equation}

As shown in the essay
\href{/physics/relativity/special-relativity/composition-of-velocities}{Composition of Velocities},
the two frames measure relative velocities of equal magnitude and opposite direction:

\begin{equation}
  \left.\vb{u}\right|_{\vb{u}'=0}=\vb{v},
  \qquad
  \left.\vb{u}'\right|_{\vb{u}=0}=-\vb{v},
  \qquad
  |\vb{u}|=|\vb{u}'|=v.
  \label{eq:reciprocal_relative_velocities}
\end{equation}

These equations provide the kinematic relations needed to compare the elapsed times measured
along the twins' different journeys.

% ============================================================
\section{Description of the Experiment}
% ============================================================

We will describe the experiment in terms of physical \textit{events}. An event is a definite
physical occurrence, such as the departure of the rocket or its arrival at planet $G$. The event
itself is independent of the reference frame, although different reference frames generally assign
different space and time coordinates to it. Consequently, a distance, an elapsed time, or a
statement of simultaneity is meaningful only after the events used in the comparison have been
specified. \par

Let us consider the Earth as the origin of the spatial coordinates of a reference frame
with coordinates $(x,y,z)$, and let us consider a planet $G$ located at a distance $l$
along the $x$-axis of this frame. \par

We assume that one of the twins remains at the origin of the Earth reference frame and
is equipped with a clock that he uses to measure time $t$. This twin does not change his
position with respect to this reference frame and measures spatial distances using the
Earth coordinate system. \par

We now consider the other twin, who will travel toward the planet $G$. We assume that,
at the beginning of the journey, the two twins occupy exactly the same position. At this
instant, the twin on the rocket introduces a spatial coordinate system $(x',y',z')$ whose
origin and axes coincide with those of the Earth frame. The time measured by the clock
carried by this twin will be denoted by $t'$. \par

Before departure, the twins synchronize their clocks and denote the departure event by $(P_0)$,
choosing it to correspond to

\[
  t=t'=0.
\]

At this instant their relative velocity is still zero, and the two spatial coordinate
systems coincide. They therefore both agree that the planet $G$ is located at a distance
$l$ from Earth along the positive $x$-axis. \par

The twin on the rocket now begins his journey toward the planet $G$, following a straight
trajectory along the direction of the common $x$-axis.

\begin{figure}[h]
  \centering
  \includegraphics[width=0.6\textwidth]{twinsBegin.png}
  \caption{Departure event}
  \label{fig:twin_departure}
\end{figure}
% ============================================================
\subsection{Acceleration to $\vb v$}
% ============================================================

We now consider the initial portion of the journey during which the rocket accelerates
until it reaches the cruise velocity $\vb v.$ We use the following two events to delimit
this phase:


\begin{itemize}
  \item  $(P_0)$ The rocket leaves Earth and begins its initial acceleration.

  \item  $(B)$ The rocket reaches the cruise velocity $\vb v$.
\end{itemize}

The twin on Earth assigns to event $(P_0)$ the coordinates

\[
  (x,y,z,t)=(0,0,0,0),
\]

and to event $(B)$ the coordinates

\[
  (x,y,z,t)=(x_1,0,0,\epsilon),
\]

where $\epsilon$ is the time required, according to the Earth clock, for the rocket
to reach its cruise velocity, and $x_1$ is the corresponding position of the rocket
relative to Earth. \par

The clock carried by the rocket is also set to zero at event $(P_0)$. Let $\epsilon'$
denote the elapsed time measured by this clock between events $(P_0)$ and $(B)$. Since
the rocket is accelerating during this interval, there is no single Lorentz transformation
that relates the Earth frame to the rocket over the entire acceleration phase. Instead,
at every instant we consider the Lorentz transformation whose velocity parameter is equal
to the instantaneous velocity of the rocket. \par

For a Lorentz transformation with velocity parameter $v$ along the $x$-axis, the
differential of the time coordinate is

\begin{equation}
  dt'
  =
  \gamma(v)
  \left(
    dt-\frac{v}{c^2}dx
  \right).
  \label{eq:Lorentz_transformation_differential_time}
\end{equation}

Looking at the differential above, we may wonder why no term proportional to $dv$
appears, since the velocity of the rocket is itself changing with time. The reason is
that the velocity $v$ appearing in the Lorentz transformation is a parameter of the
transformation. For each Lorentz transformation it is therefore held fixed when the
differential with respect to the spacetime coordinates is taken. Thus,

\[
  dt'
  =
  \gamma(v)
  \left(
    dt-\frac{v}{c^2}dx
  \right),
\]

contains no contribution proportional to $dv$. \par

During the acceleration phase, the value of the parameter $v$ changes along the
trajectory of the rocket. We therefore use, at each event of the trajectory, the
Lorentz transformation whose velocity parameter is equal to the velocity of the rocket
at that event. In this sense, $dt'$ represents the infinitesimal time measured in the
corresponding instantaneous rocket frame; it should not be interpreted as the
differential of a single inertial time coordinate $t'$ extending over the entire
acceleration phase. Rather, we consider a family of Lorentz
transformations

\[
  \boldsymbol{\Lambda}\bigl(v(t)\bigr),
\]

and at each instant we use the member of this family whose fixed parameter is equal to
the instantaneous velocity of the rocket. The parameter labeling the transformation
changes along the trajectory, but it does not contribute to the differential within
each individual transformation.  \par
Along the trajectory of the rocket,

\[
  dx=v(t)\,dt,
\]

and therefore \eqref{eq:Lorentz_transformation_differential_time} becomes

\begin{equation}
  dt'
  =
  \gamma(v)
  \left(
    1-\frac{v^2}{c^2}
  \right)dt
  =
  \frac{dt}{\gamma(v)}.
  \label{eq:rocket_clock_differential}
\end{equation}

The elapsed time measured by the clock carried by the rocket during the acceleration
phase is consequently

\begin{equation}
  \epsilon'
  =
  \int_0^\epsilon
  \frac{dt}{\gamma(v(t))}
  =
  \int_0^\epsilon
  \sqrt{1-\frac{v^2(t)}{c^2}}\,dt.
  \label{eq:acceleration_elapsed_time}
\end{equation}

The precise relation between $\epsilon$ and $\epsilon'$ therefore depends on the
acceleration profile $v(t)$ of the rocket. No assumption is made here that the velocity
is small compared with the speed of light; on the contrary, the cruise velocity may
be arbitrarily close to $c$. \par


We will assume that the acceleration phase is very short compared with the duration
of the complete journey. More precisely, if $T$ and $T'$ denote characteristic time
scales of the journey as measured in the Earth frame and by the clock carried by the
rocket, respectively, we assume that

\begin{equation}
  \epsilon \ll T,
  \qquad
  \epsilon' \ll T'.
  \label{eq:negligible_acceleration_times}
\end{equation}

The times spent accelerating are therefore negligible compared with the times spent
traveling at cruise velocity. In the remainder of the analysis we will neglect both
$\epsilon$ and $\epsilon'$, while allowing the cruise velocity $\vb v$ to remain fully
relativistic. \par

We further assume that the distance $x_1$ traveled by the rocket during the initial
acceleration phase is very small compared with the distance $l$ between Earth and
planet $G$,

\[
  x_1 \ll l.
\]

We will therefore also neglect $x_1$ in the remainder of the analysis.

%%%%%%%%%%%%%%%%%%%%%%%%%%%%%%%%%%%%%%%%%%%
 
% ============================================================
\subsection{Trip to $G$}
% ============================================================

We now consider the journey of the rocket to the planet $G$. Since the motion takes
place entirely along the direction of the common $x$-axis, we will omit the coordinates
$y$ and $z$ in what follows. \par


We consider the following two events:

\begin{itemize}
  \item $(P_0)$ The rocket departs from Earth.

  \item $(P_1)$ The rocket reaches the planet $G$ with velocity $\vb v$ and begins
  to decelerate.
\end{itemize}

In light of the approximations introduced for the initial acceleration phase, we neglect
both the time and the distance between events $(P_0)$ and $(B)$. Thus, for the cruise-phase
calculation, event $(B)$ is treated as coincident with $(P_0)$ to the accuracy of our
approximation. Accordingly, both twins assign to event $(P_0)$ the coordinates

\[
  (x,t)=(0,0),
  \qquad
  (x',t')=(0,0).
\]

The twin on Earth assigns to event $(P_1)$ the coordinates

\[
  (x,t)=(l,t_1),
\]

where $t_1$ is the time required for the rocket to travel from Earth to planet $G$.
Since the rocket moves with constant speed $v=|\vb v|$ during this portion of the
journey,

\begin{equation}
  t_1=\frac{l}{v}.
  \label{eq:outbound_time_earth}
\end{equation}

During the outbound cruise, let $F'_{\mathrm{out}}$ denote the inertial frame comoving with
the rocket. The rocket moves with velocity $+\vb v$ relative to the Earth frame $F$, while
Earth and planet $G$ move with velocity $-\vb v$ relative to $F'_{\mathrm{out}}$. The twin on
the rocket assigns to event $(P_1)$ the coordinates

\[
  (x',t')=(0,t_1'),
\]

where $t_1'$ is the elapsed time measured by his clock during the journey to planet
$G$. There are two different ways in which he can determine $t_1'$. \par

The first is to use the Lorentz transformation directly. Applying it to event $(P_1)$,
whose coordinates in the Earth frame are

\[
  (x,t)=\left(l,\frac{l}{v}\right),
\]

we obtain

\begin{equation}
  \begin{aligned}
    x'
    &=
    \gamma(x-vt)
    =
    \gamma\left(
      l-v\frac{l}{v}
    \right)
    =
    0,
    \\
    t_1'
    &=
    \gamma\left(
      t-\frac{v}{c^2}x
    \right)
    \\
    &=
    \gamma\left(
      \frac{l}{v}-\frac{vl}{c^2}
    \right)
    =
    \frac{1}{\gamma}\frac{l}{v}.
  \end{aligned}
  \label{eq:outbound_time_rocket_lorentz}
\end{equation}

Therefore,

\begin{equation}
  t_1'=\frac{t_1}{\gamma}.
  \label{eq:outbound_time_relation}
\end{equation}

There is, however, another way for the twin on the rocket to understand this result.
As soon as the rocket has reached its cruise velocity $\vb v$, the distance between
Earth and planet $G$, measured in the rocket frame, is smaller than the corresponding
distance $l$ measured in the Earth frame. \par

To understand why, let us first consider how the twin on the rocket measures a length
belonging to an object at rest in the Earth frame. As an example, consider the diameter
of the Earth along the direction of motion. Let its diameter measured in the Earth frame
be

\[
  d=2r.
\]

A measurement of this diameter in the rocket frame requires the positions of the two
opposite endpoints of the diameter to be determined \textit{simultaneously in the
rocket frame}. We therefore introduce two events:

\begin{itemize}
  \item $(E_+)$ The position of the endpoint of the Earth on the positive side of the
  $x$-axis is measured.

  \item $(E_-)$ The position of the endpoint of the Earth on the negative side of the
  $x$-axis is measured.
\end{itemize}

Since these two measurements are simultaneous in the rocket frame, their coordinates
can be written as

\begin{equation}
  \begin{aligned}
    E_+ &: \qquad (x',t')=(x_+',t_+'), \\
    E_- &: \qquad (x',t')=(x_-',t_-').
  \end{aligned}
  \label{eq:earth_diameter_events_rocket}
\end{equation}

The diameter measured by the twin on the rocket is therefore

\begin{equation}
  d'=x_+'-x_-'.
  \label{eq:earth_diameter_rocket_definition}
\end{equation}

It is important to emphasize that the measurement of a spatial distance in a given
reference frame requires the positions of its endpoints to be measured at the same
time in that frame. The events $E_+$ and $E_-$ are simultaneous in the rocket frame,
but they are not simultaneous in the Earth frame. \par

Since the Earth is at rest in the Earth frame and centered at the origin, the same two
events have coordinates

\begin{equation}
  \begin{aligned}
    E_+ &: \qquad (x,t)=(r,t_+), \\
    E_- &: \qquad (x,t)=(-r,t_-).
  \end{aligned}
  \label{eq:earth_diameter_events_earth}
\end{equation}

Because the two events are simultaneous in the rocket frame,

\[
  t_+'-t_-'=0,
\]

using the Lorentz transformation for the time coordinate, we obtain

\begin{equation}
  \begin{aligned}
    0
    &=
    \gamma\left(
      t_+-\frac{v}{c^2}r
    \right)
    -
    \gamma\left(
      t_- -\frac{v}{c^2}(-r)
    \right)
    \\
    &=
    \gamma\left[
      t_+-t_-
      -
      \frac{v}{c^2}(2r)
    \right].
  \end{aligned}
\end{equation}

Since $d=2r$, it follows that

\begin{equation}
  t_+-t_-=\frac{v}{c^2}d.
  \label{eq:earth_diameter_relativity_simultaneity}
\end{equation}

Thus, although $E_+$ and $E_-$ occur simultaneously for the twin on the rocket,
event $E_+$ occurs after event $E_-$ according to the twin on Earth. This is a direct
manifestation of the relativity of simultaneity. \par

We can now use the spatial part of the Lorentz transformation to determine the diameter
measured in the rocket frame. Taking the difference between the transformed positions
of the two events gives

\begin{equation}
  \begin{aligned}
    d'
    &=
    x_+'-x_-'
    \\
    &=
    \gamma\left[
      (r-(-r))-v(t_+-t_-)
    \right].
  \end{aligned}
\end{equation}

Using \eqref{eq:earth_diameter_relativity_simultaneity}, we obtain

\begin{equation}
  \begin{aligned}
    d'
    &=
    \gamma\left(
      d-\frac{v^2}{c^2}d
    \right)
    \\
    &=
    \frac{d}{\gamma}.
  \end{aligned}
  \label{eq:length_contraction_example}
\end{equation}

The twin on the rocket therefore measures a smaller longitudinal diameter for the
Earth than the twin at rest with respect to it. More generally, the separation between
two points that are at rest in the Earth frame and separated along the direction of
relative motion is measured in the rocket frame to be


\begin{equation}
  l'=\frac{l}{\gamma}.
  \label{eq:earth_planet_distance_rocket}
\end{equation}


In particular, Earth and planet $G$ are both at rest in the Earth frame and are
separated by the distance $l$. The twin on the rocket therefore measures their
separation to be $l'$. \par

The following figure summarizes the outbound journey in both reference frames. In the Earth
frame, Earth and planet $G$ are at rest and separated by the distance $l$, while the
rocket moves between them with speed $v$. In the rocket frame, the rocket is at rest,
while Earth and planet $G$ move with the same speed $v$ in the opposite direction, but
their separation is reduced to $l/\gamma$. 

In the rocket-frame diagram, the rocket is shown midway between Earth and planet $G$
only for visual convenience. This midpoint has no special physical significance. At the
instant represented in the figure,

\[
  x'_{\mathrm E}=-\frac{l'}{2},
  \qquad
  x'_G=\frac{l'}{2},
\]

with

\[
  l'=\frac{l}{\gamma}.
\]

At any other instant during the cruise, the individual coordinates of Earth and planet
$G$ would be different, while their simultaneous separation in the rocket frame would
remain equal to $l'$.

\begin{figure}[h]
  \centering
  \includegraphics[width=0.6\textwidth]{rocketOutbound.png}
  \caption{Outbound trip: Earth frame $F$ versus the outbound rocket frame $F'_{\mathrm{out}}$.}
  \label{fig:twin_outbound_frames}
\end{figure}


As shown previously in \eqref{eq:reciprocal_relative_velocities}, if the rocket moves
with velocity $\vb v$ relative to Earth, then Earth and planet $G$ move with velocity
$-\vb v$ relative to the rocket. The magnitude of the relative velocity is therefore
the same in both descriptions and is equal to $v$. Consequently, according to the
rocket twin, planet $G$ approaches him with speed $v$, and the time required for the
planet to reach him is

\begin{equation}
  t_1'
  =
  \frac{l'}{v}
  =
  \frac{1}{\gamma}\frac{l}{v}.
  \label{eq:outbound_time_rocket_length}
\end{equation}

This is identical to the result obtained directly from the Lorentz transformation in
\eqref{eq:outbound_time_rocket_lorentz}. Since $\gamma>1$ for $v>0$, we also have

\begin{equation}
  t_1'<t_1.
\end{equation}

Thus, the two twins agree on the events marking the beginning and the end of this
portion of the journey, but they assign different spatial distances and different
elapsed times to it. For the twin on the rocket, the distance between Earth and planet
$G$ is shorter, and correspondingly less time elapses before the rocket and the planet
meet.

%%%%%%%%%%%%%%%%%%%%%%%%%%%%%%%%%%%%%%%%%%%
 
% ============================================================
\subsection{Trip Back to Earth}
% ============================================================

Once the twin on the rocket reaches planet $G$, he reverses his motion and begins the
journey back to Earth. As in the initial acceleration phase, we assume that the time
required to reverse the motion is very small compared with the duration of the journey,
and we neglect both this time interval and the corresponding displacement. \par

We denote by $(P_2)$ the event at which the traveling twin reunites with the twin on Earth.
For the return cruise, let $F'_{\mathrm{in}}$ denote the inertial frame comoving with the
rocket. The rocket moves with velocity $-\vb v$ relative to the Earth frame $F$, while Earth
and planet $G$ move with velocity $+\vb v$ relative to $F'_{\mathrm{in}}$. The magnitude of
the relative velocity is therefore again $v$. Moreover, for the same reason discussed in
the previous section, the distance between Earth and planet $G$ measured in the inbound
rocket frame is

\[
  l'=\frac{l}{\gamma}.
\]

Consequently, the elapsed time measured by the twin on the rocket during the return
journey is

\begin{equation}
  t_2'
  =
  \frac{l'}{v}
  =
  \frac{1}{\gamma}\frac{l}{v}.
  \label{eq:return_time_rocket}
\end{equation}

For the twin on Earth, on the other hand, the rocket travels the distance $l$ from
planet $G$ back to Earth with speed $v$. The elapsed time for the return journey is
therefore

\begin{equation}
  t_2
  =
  \frac{l}{v}.
  \label{eq:return_time_earth}
\end{equation}

The outbound and return portions of the journey have the same duration in each
corresponding description,

\[
  t_2=t_1,
  \qquad
  t_2'=t_1'.
\]

Within the approximation that the launch and turnaround durations and displacements are
negligible compared with the two cruise portions, the total elapsed time measured by the
clock carried by the rocket is therefore

\begin{equation}
  \begin{aligned}
    T'
    &=
    t_1'+t_2'
    \\
    &=
    2\frac{1}{\gamma}\frac{l}{v},
  \end{aligned}
  \label{eq:total_journey_time_rocket}
\end{equation}

whereas the total elapsed time measured by the clock on Earth is

\begin{equation}
  \begin{aligned}
    T
    &=
    t_1+t_2
    \\
    &=
    2\frac{l}{v}.
  \end{aligned}
  \label{eq:total_journey_time_earth}
\end{equation}

Since $\gamma>1$ for any nonzero relative velocity,

\begin{equation}
  T'<T.
  \label{eq:twin_elapsed_time_difference}
\end{equation}

Thus, at the reunion event $(P_2)$ on Earth, the clock carried by the traveling twin
records a smaller elapsed time than the clock that remained on Earth. Since each twin's
clock measures the time experienced by that twin, the traveling twin is younger when
the two are reunited. The quantity $T'$ is the elapsed time recorded by the same continuous
onboard clock throughout the entire journey; it is not a coordinate time belonging to one
global inertial frame $F'$.
 
The complete journey is summarized in the following figure. The upper row
represents the motion in the Earth frame, while the lower row represents the corresponding
description during the constant-velocity portions of the rocket's motion. The short
turnaround interval is shown explicitly but is assumed to be negligible compared with the
duration of the two cruise portions. \par

\begin{figure}[h]
  \centering
  \includegraphics[width=0.6\textwidth]{rocketInbound.png}
  \caption{Complete journey in the Earth frame $F$ and in the corresponding outbound and
  inbound rocket frames $F'_{\mathrm{out}}$ and $F'_{\mathrm{in}}$.}
  \label{fig:twin_complete_journey}
\end{figure}



As in the previous figure, the precise positions shown during the cruise are schematic.
What is physically relevant is that Earth and planet $G$ are separated by the distance
$l$ in the Earth frame and by the contracted distance $l/\gamma$ in the corresponding
outbound or inbound rocket frame.
% ============================================================
\subsection{Events and the Apparent Asymmetry of Time Dilation}
% ============================================================

We will now return to the question raised at the beginning of the essay and address
the apparent asymmetry in the motion of the two twins. One of the apparent
contradictions comes from the familiar statement that a clock observed from a
reference frame in which it is moving appears to run more slowly. It therefore seems
reasonable that, when the rocket reaches planet $G$, the clock carried by the rocket
has recorded less time than the clock on Earth. By symmetry, however, we might also
expect the twin on the rocket to find that the clock on Earth has recorded less time
than his own. \par

The resolution of this apparent contradiction requires us to be very precise about
the events being compared. The familiar formulas for time dilation do not apply to
arbitrary events: they apply to particular pairs of events selected by specific
conditions. \par

Let us first consider the arrival of the rocket at planet $G$ from the point of view
of the twin on Earth. The arrival event is

\[
  P_1:
  \qquad
  (x,t)
  =
  \left(
    l,\frac{l}{v}
  \right).
\]

At the same time in the Earth frame, the twin on Earth is still at the origin of $F$.
We therefore consider the event

\[
  Q_1:
  \qquad
  (x,t)
  =
  \left(
    0,\frac{l}{v}
  \right).
\]

Thus, in the Earth frame, the two events

\[
  P_1
  \qquad\text{and}\qquad
  Q_1
\]

are simultaneous. The first occurs at planet $G$, while the second occurs at the
position of the twin on Earth. In this description, the Earth twin is at rest and his
position remains fixed at $x=0$. \par

Let us now describe the same arrival event $P_1$ from the outbound rocket frame
$F'_{\mathrm{out}}$. Its coordinates are

\[
  P_1:
  \qquad
  (x',t')
  =
  \left(
    0,\frac{l}{\gamma v}
  \right).
\]

At this same time in $F'_{\mathrm{out}}$, the twin on Earth is not at a fixed spatial
position. He is moving toward the negative $x'$ direction with velocity $-\vb{v}$.
Since the distance between Earth and planet $G$ measured in the rocket frame is

\[
  l'=\frac{l}{\gamma},
\]

at the instant at which the rocket reaches planet $G$, the Earth twin is located at

\[
  x'=-\frac{l}{\gamma}.
\]

We therefore introduce the event

\[
  R_1:
  \qquad
  (x',t')
  =
  \left(
    -\frac{l}{\gamma},
    \frac{l}{\gamma v}
  \right).
\]

Thus, in the rocket frame, the two events

\[
  P_1
  \qquad\text{and}\qquad
  R_1
\]

are simultaneous. Again, one event is the arrival of the rocket at planet $G$, while
the other occurs at the position of the twin on Earth at that same time according to
the rocket frame. The difference is that, in this description, the Earth twin is
moving rather than remaining at a fixed spatial coordinate. \par

We can now determine how the event $R_1$ is recorded in the Earth frame by using the
inverse Lorentz transformation,

\[
  x=\gamma(x'+vt'),
  \qquad
  t=\gamma\left(t'+\frac{v}{c^2}x'\right).
\]

For the spatial coordinate we obtain

\begin{align}
  x(R_1)
  &=
  \gamma
  \left(
    -\frac{l}{\gamma}
    +
    v\frac{l}{\gamma v}
  \right)
  \nonumber\\
  &=
  0,
\end{align}

as expected, since $R_1$ is the position of the twin on Earth. For the
time coordinate,

\begin{align}
  t(R_1)
  &=
  \gamma
  \left(
    \frac{l}{\gamma v}
    -
    \frac{v}{c^2}\frac{l}{\gamma}
  \right)
  \nonumber\\
  &=
  \frac{l}{v}
  \left(
    1-\frac{v^2}{c^2}
  \right)
  \nonumber\\
  &=
  \frac{l}{\gamma^2 v}.
\end{align}

Since the elapsed time measured by the rocket clock between departure and $P_1$ is

\[
  t_1'=\frac{l}{\gamma v},
\]

we may equivalently write

\[
  t(R_1)=\frac{t_1'}{\gamma}.
\]

We therefore recover the familiar reciprocal time-dilation result: according to the
outbound rocket frame, the clock on Earth has advanced by less than the clock carried
by the rocket. \par

There is nevertheless no contradiction with the result obtained previously. The
Earth event $Q_1$ used in the Earth-frame description and the Earth event $R_1$ used
in the rocket-frame description are not the same event. Indeed,

\[
  t(Q_1)=\frac{l}{v},
  \qquad
  t(R_1)=\frac{l}{\gamma^2 v}.
\]

Thus,

\[
  Q_1\neq R_1.
\]

The two events are selected by different simultaneity conditions. In the Earth frame,

\[
  t(P_1)=t(Q_1),
\]

while in the outbound rocket frame,

\[
  t'(P_1)=t'(R_1).
\]

The two reference frames therefore associate different events on the Earth twin's
reference system with the same arrival event $P_1$. \par

This distinction is the key point. The arrival event $P_1$ is not arbitrary: it is
constrained by the physical setup of the experiment. Planet $G$ remains at the fixed
position $x=l$ in the Earth frame, while the rocket travels with speed $v$. Therefore
the rocket can meet planet $G$ only when

\[
  l=vt,
\]

which fixes the Earth-frame time of the arrival event to be

\[
  t(P_1)=\frac{l}{v}.
\]

The event selected by the reciprocal time-dilation comparison in
$F'_{\mathrm{out}}$ is instead $R_1$, a different event on the Earth reference system. The
two apparently contradictory statements therefore compare different pairs of events.
\par


\begin{figure}[h]
  \centering
  \includegraphics[width=0.6\textwidth]{eventsRole.png}
  \caption{The role of events in the comparison of elapsed times.}
  \label{fig:role_of_events}
\end{figure}

This illustrates a principle that is fundamental throughout special relativity:
before applying a formula for time dilation, length contraction, or simultaneity, we
must first identify precisely the events to which that formula refers.

\end{document}
